\section[Analyse Intel]{Analyse des résultats sur Intel}

\begin{frame}{Intel Fortran MPI : duel \texttt{ifx} / \texttt{gfortran}}
\small

\begin{columns}[T,onlytextwidth]
    \begin{column}{0.44\textwidth}
        \begin{block}{Contexte}
            \begin{itemize}\setlength{\itemsep}{1em}
                \item Comparaison sur la version \textbf{Fortran MPI officielle} de NPB MG.
                \item Évaluation de l'efficacité du compilateur Intel face au standard \texttt{gfortran}.
            \end{itemize}
        \end{block}
    \end{column}
    \hfill 
    \begin{column}{0.52\textwidth}
        \begin{block}{Métriques principales}
            \centering
            \renewcommand{\arraystretch}{1.1}
            \footnotesize
            \begin{tabular}{|l|c|c|}
                \hline
                \textbf{Métrique} & \textbf{ifx} & \textbf{gfortran} \\ \hline
                Total Time (s)      & 18,10 & 22,92 \\ \hline
                Max Active (s)      & 17,69 & 22,53 \\ \hline
                Avg Active (s)      & 17,60 & 22,40 \\ \hline
                Activity Ratio (\%) & 97,3  & 97,8  \\ \hline
                Threads actifs      & 7,779 & 7,819 \\ \hline
                Scalability Gap     & 1,00  & 1,27  \\ \hline
            \end{tabular}
        \end{block}
    \end{column}
\end{columns}
\vfill 
\begin{exampleblock}{Analyse du gain}
    \centering
    \small
    L'utilisation de \texttt{ifx} permet de passer de \textbf{22,92 s} à \textbf{18,10 s}, soit un gain de performance de \textbf{21 \%} directement visible sur le temps total.
\end{exampleblock}

\end{frame}

\begin{frame}{Intel Fortran MPI : stabilité et scalabilité}
\footnotesize 

\begin{columns}[T,onlytextwidth]
    \begin{column}{0.48\textwidth}
        \begin{block}{\texttt{-S1} : stabilité (10 runs)}
            \begin{itemize}\setlength{\itemsep}{0pt} 
                \item Total Time : \textbf{18,09 s}
                \item Max Active moyen : \textbf{17,60 s}
                \item Min / Max : \textbf{16,46 s} / \textbf{17,88 s}
                \item Médiane : \textbf{17,71 s}
                \item \textbf{0 outlier}
            \end{itemize}
        \end{block}
    \end{column}

    \begin{column}{0.48\textwidth}
        \begin{block}{\texttt{-WS} : 1 à 8 processus}
            \centering
            \setlength{\tabcolsep}{4pt}
            \renewcommand{\arraystretch}{0.9} 
            \begin{tabular}{|c|c|c|}
                \hline
                \textbf{Config} & \textbf{Temps (s)} & \textbf{Gap} \\ \hline
                1 proc & 21,18 & 1,00 \\ \hline
                2 proc & 16,03 & 1,51 \\ \hline
                4 proc & 15,02 & 2,84 \\ \hline
                8 proc & 14,30 & 5,40 \\ \hline
            \end{tabular}
        \end{block}
    \end{column}
\end{columns}

\begin{block}{Constats}
    \begin{itemize}\setlength{\itemsep}{0pt}
        \item \textbf{Stabilité} : Comportement reproductible validant le protocole.
        \item \textbf{Scalabilité} : Gain limité par la saturation mémoire (\textit{Memory Bound}).
    \end{itemize}
\end{block}

\end{frame}

\begin{frame}{Identification des sections critiques (Option -R1)}
    \begin{columns}[T]
        \begin{column}{0.45\textwidth}
            \begin{itemize}
                \fontsize{8}{9}\selectfont
                \item Localisation des fonctions consommant le plus de temps CPU.
                \item \textbf{resid} est la routine dominante (35,48 \%), suivie par des temps de synchronisation importants.
                \item Profil analysé : \textbf{NPB-CPP (C++ OpenMP)} compilé avec \texttt{icpx}.
                \item Les noyaux stencil restent dominants, mais une part importante du temps est perdue dans le runtime OpenMP.
            \end{itemize}
            \vspace{0.1cm}
            \begin{table}
                \fontsize{6}{7}\selectfont
                \centering
                \begin{tabular}{|l|c|}
                    \hline
                    \textbf{Fonction} & \textbf{\% Temps CPU} \\ \hline
                    \texttt{resid}    & 35,48 \% \\ \hline
                    \texttt{kmp\_flag} (Sync) & 25,24 \% \\ \hline
                    \texttt{zran3}    & 17,23 \% \\ \hline
                    \texttt{psinv}    & 15,65 \% \\ \hline
                \end{tabular}
            \end{table}
        \end{column}
        \begin{column}{0.55\textwidth}
            \vspace{-0.4cm}
            \inputImage{figures/top_functions_intel_r1.png}{Analyse -R1 : Routines dominantes (Intel)}{0.9\textwidth}{!}{fig:top_functions_intel_r1}
        \end{column}
    \end{columns}
\end{frame}

\begin{frame}{Intel C++ OpenMP : duel de compilateurs (\texttt{icpx} vs \texttt{clang++})}

    \begin{itemize}\setlength{\itemsep}{0pt}
        \fontsize{8}{9}\selectfont
        \item Comparaison sur le même code C++ OpenMP : seule la \textbf{toolchain} change.
        \item Objectif : isoler l'impact du compilateur sur la performance brute et le runtime.
    \end{itemize}

    \centering
    \inputImage{figures/comparison_duel_report_intel.png}{Comparaison Intel : icpx vs clang++}{0.7\textwidth}{!}{fig:comparison_duel_report_intel}

    \vspace{-0.8cm}
    \begin{columns}[T,onlytextwidth]
        \begin{column}{0.48\textwidth}
            \begin{exampleblock}{\footnotesize \texttt{icpx} (r0)}
                \fontsize{7}{8}\selectfont
                \begin{itemize}\setlength{\itemsep}{0pt}
                    \item Temps : \textbf{23,83 s}.
                    \item Array access : \textbf{97,4 \%}.
                    \item Activity Ratio : \textbf{72,8 \%}.
                \end{itemize}
            \end{exampleblock}
        \end{column}
        \begin{column}{0.48\textwidth}
            \begin{alertblock}{\footnotesize \texttt{clang++} (r1)}
                \fontsize{7}{8}\selectfont
                \begin{itemize}\setlength{\itemsep}{0pt}
                    \item Temps : \textbf{22,55 s}.
                    \item Array access : \textbf{95,9 \%}.
                    \item Activity Ratio : \textbf{64,4 \%}.
                \end{itemize}
            \end{alertblock}
        \end{column}
    \end{columns}
\end{frame}

\begin{frame}{Intel C++ OpenMP : stabilité et scalabilité}
    \begin{columns}[T,onlytextwidth]
        \begin{column}{0.46\textwidth}
            \inputImage{figures/stability_analysis_intel_s1.png}
                       {Stabilité -S1 (Intel)}
                       {1.0\linewidth}
                       {!}   
                       {fig:stability_analysis_intel_s1}
            
            \vspace{-0.2cm}
            \begin{itemize}\setlength{\itemsep}{2pt}
                \fontsize{7}{8}\selectfont
                \item Temps moyen : \textbf{23,68 s}.
                \item Min -- Max : \textbf{23,39 s} -- \textbf{23,82 s}.
                \item \textbf{0 outlier} : mesure fiable.
            \end{itemize}
        \end{column}

        \hfill 

        \begin{column}{0.46\textwidth}
            \inputImage{figures/scalability_speedup_intel_ws.png}
                       {Scalabilité -WS (Intel)}
                       {1.0\linewidth} 
                       {!}           
                       {fig:scalability_speedup_intel_ws}
            
            \vspace{-0.2cm}
            \begin{itemize}\setlength{\itemsep}{2pt}
                \fontsize{7}{8}\selectfont
                \item Séquentiel (r0) : \textbf{25,91 s}.
                \item OMP 4 (r2) : \textbf{23,33 s} (optimum).
                \item OMP 8 (r3) : \textbf{23,78 s} (saturation).
            \end{itemize}
        \end{column}
    \end{columns}
\end{frame}

\begin{frame}{Synthèse Intel : Fortran MPI et C++ OpenMP}
    \begin{itemize}
        \fontsize{8}{9}\selectfont
        \item \textbf{Fortran MPI} : le compilateur Intel \texttt{ifx} est le résultat le plus convaincant sur la plateforme Intel.
        \begin{itemize}
            \fontsize{8}{9}\selectfont
            \item \texttt{ifx} réduit le temps de \textbf{22,92 s} à \textbf{18,10 s} face à \texttt{gfortran}.
            \item La stabilité est bonne et la scalabilité reste utile jusqu'à \textbf{8 processus}.
        \end{itemize}
        \item \textbf{C++ OpenMP} : le code reste limité par la structure du runtime et la synchronisation.
        \begin{itemize}
            \fontsize{8}{9}\selectfont
            \item \texttt{kmp\_wait} représente \textbf{25,24 \%} du temps.
            \item \texttt{icpx} reste légèrement derrière \texttt{clang++} en temps total.
            \item Le meilleur point est obtenu à \textbf{4 threads}, puis la montée en charge plafonne.
        \end{itemize}
    \end{itemize}

    \begin{block}{Conclusion}
        \fontsize{8}{9}\selectfont
        Sur la plateforme Intel, \textbf{le Fortran MPI profite bien du compilateur Intel}, tandis que le \textbf{C++ OpenMP reste surtout limité par l'efficacité du parallélisme partagé}. Les pistes prioritaires sont donc différentes selon le paradigme.
    \end{block}
\end{frame}
